\section{Conditional variance at the exact cap scale}\label{sec:variance}

The observation removes the part of a test that depends on the leading
spectral coordinates. We now show that its remaining variance costs at most
$Cp^{-2}$ times the original heat-bath energy. The exponent $2$ is essential:
integrating a covariance bound $(1+\varepsilon)p^{-2}$ with a fixed
$\varepsilon>0$ would instead produce a power loss. The integrated estimates
of Section~\ref{sec:observations} allow us to replace this fixed error by
one that is integrable in logarithmic scale.

\begin{proposition}[Conditional variance]\label{prop:variance-main}
There is a fixed $p_*>0$ with the following property. For every
$\varepsilon>0$ there are finite constants $C$, $\xi>0$ and disorder events
$\mathcal V_n$ with $\liminf_n\mathbb P(\mathcal V_n)\ge1-\varepsilon$
such that, on $\mathcal V_n$, simultaneously for every real function $f$
on the cube and every $r\le p\le p_*$,
\begin{equation}\label{eq:source-16}
 V_p(f):=Q_\mu[\Var(f\mid h_p)]
 \le Cp^{-2}\cE(f)+Ce^{-n^\xi}\|f\|_\infty^2.
\end{equation}
The constants and the disorder event are independent of $f$. The inequality
also holds componentwise and may be summed for vector-valued functions.
\end{proposition}

In particular, the remainder is negligible for any fixed polynomial bound
on the supremum norm, including for disorder-dependent tests. We first
prove the requisite covariance improvement, then establish a variance bound
at a fixed upper scale, and finally integrate the posterior variance identity.

\subsection{A sharper upper covariance bound}

We work initially on the disorder event $\mathcal G_n$ of
Section~\ref{sec:observations}. Its fixed tolerances will be chosen below;
no tolerance depends on $n$. Probabilities in the next proposition are
measured under $\mathbf P$ before restriction to this event.

\begin{proposition}[Positive covariance error]\label{prop:variance-curvature}
There are fixed $a,\xi,C>0$ and a sufficiently small fixed $p_*>0$ such
that, for each deterministic $p\in[r,p_*]$ away from the activation
scales,
\begin{equation}\label{eq:source-17}
 p^2K_p\preceq\bigl(1+C(p^a+k_p^{-a})\bigr)\Id
\end{equation}
outside an event $\mathcal N_{n,p}$ satisfying
$\mathbf P(\mathcal N_{n,p}\cap\mathcal G_n)\le Ce^{-n^\xi}$.
The exceptional event may depend on $p$; its probability bound has common
constants for the entire scale range. By Fubini's theorem the same bound
holds after integration in $d\log p$, with a possibly smaller $\xi$.
\end{proposition}

The gain comes from comparing a small cap to a larger one. The radial
covariance deficit permits us to put all positive errors in the tangent
space of the parent field. Their restriction to the much smaller child
space is then controlled by the unused Haar rotation.

Fix a small constant $\gamma>0$, a child scale $c$, and a parent scale $p$
such that $x=c/p\in[\gamma,2\gamma]$. Put $\Delta=B_p-B_c$.
Conditional on $h_c$, the parent field has the transition law
\[
 d\mathsf T_{c,p,h_c}(y)
 =\frac{e^{\phi_p(y)}\,dN(h_c,\Delta)(y)}
        {\int e^{\phi_p(z)}\,dN(h_c,\Delta)(z)}.
\]
The Gaussian is taken on the positive-increment space. The observation
precision increment, compressed to the child space, is
$\Delta_c=(P_c\Delta P_c)|_{\mathcal H_c}=(1-x^2)p^2\Id$.
Writing $[M]_{\mathcal H_c}$ for this compression of any operator $M$,
differentiation of the Gaussian convolution gives
\begin{equation}\label{eq:source-18}
 K_c=\Delta_c^{-1}
       [\Cov_{\mathsf T_{c,p,h_c}}(h_p)]_{\mathcal H_c}\Delta_c^{-1}
        -\Delta_c^{-1}.
\end{equation}

\begin{lemma}[An annular extension for the transition]\label{lem:variance-extension}
Choose $\eta>0$ sufficiently small compared with $\gamma^2$, then the
coarse curvature tolerance $\varepsilon_0$ sufficiently small compared with
$\eta$. The annulus width $s_0$ and the fixed upper scale may be chosen
so that the following holds on $\mathcal G_n$. For every child field with
$|\alpha_c(h_c)|\le s_0c^2$, the parent potential has an extension
$\overline\phi_p$ with upper Hessian bound $(1+\eta)p^{-2}\Id$ that agrees
with $\phi_p$ near $|\alpha_p|\le s_0p^2/2$. Under both the true transition
and the transition obtained by using $\overline\phi_p$, the complement of
this agreement region has probability at most $Ce^{-a_1k_p}$. The same
bound holds with any fixed polynomial moment of
$\|h_p\|/(p^2M_p)$ inserted, after decreasing $a_1>0$.
\end{lemma}

\begin{proof}
We give the construction and its tail comparison, following the cap
extension method of \cite[Proposition~16.1]{CM}. On
$|\alpha_p|\le3s_0p^2/4$, take the infimum of the quadratics
\[
 \phi_p(z)+\nabla\phi_p(z)\cdot(y-z)
             +\frac{1+\eta}{2p^2}\|y-z\|^2.
\]
Each of these supports lies above $\phi_p$ on that smaller annulus.
Choose $s_0$ small compared with $\eta$ and $\gamma^2$. The radial buffer
between the smaller and larger annuli then determines a fixed
$\beta_{\mathrm{buf}}>0$ such that every segment starting in the smaller
annulus and of length at most $\beta_{\mathrm{buf}}p^2M_p$ remains in
the larger one. For such pairs $y,z$, the supporting inequality follows
by integrating the coarse Hessian bound along their segment.
For $d=y-z$ with $\|d\|\ge\beta_{\mathrm{buf}}p^2M_p$, freeze the
spherical saddle at $z$. The frozen Gaussian bound bounds the spherical
logarithmic increment by the quadratic with Hessian
$(p^2+\alpha_p(z))^{-1}\Id$, its saddle gradient, and a bounded density
error. The static value and gradient errors from
Lemma~\ref{lem:edge-static-interpolation} are at most
$\epsilon_v k_p$ and $\epsilon_gM_p$. Relative to $\|d\|^2/p^2$, their
combined contribution, including the density error, is at most
\[
 \frac{2\epsilon_v+C/k_p}{\beta_{\mathrm{buf}}^2}
       +\frac{\epsilon_g}{\beta_{\mathrm{buf}}}.
\]
The frozen Hessian is at most $(1+Cs_0)p^{-2}\Id$. Thus, after choosing
$s_0$, then $\beta_{\mathrm{buf}}$, then the static errors, the excess
quadratic term dominates this residual for all sufficiently large $n$.
This proves the supporting property,
and the infimum is a semiconcave extension agreeing on the smaller annulus.

For the transition tails freeze instead the child saddle
$\alpha=\alpha_c(h_c)$. As reference use the Gaussian transition tilted
by the parent quadratic $\|y\|^2/(2(p^2+\alpha))$ with no linear term.
Writing $R_u=R_{u,\alpha}$, its law is characterized by
\[
 R_ph_p\sim N(R_ch_c,R_c-R_p),
 \qquad R_c^{-1}=R_p^{-1}-\Delta.
\]
Since the child saddle satisfies its norm constraint, the expected squared
norm of $h_p$ is $\ell_p(\alpha)^2$. Its covariance norm is at most
$C\gamma^{-2}p^2$. Therefore deviations in radius by $vp^2M_p$ cost
\[
 C\exp(-a_0\gamma^2k_pv^2),
 \qquad v\ge C(\gamma\sqrt{k_p})^{-1}.
\]
The radius $\ell_p(\alpha)$ lies strictly inside the parent agreement
annulus, because $|\alpha|\le s_0c^2$.

Calibrate the reference quadratic to the spherical potential at this
radius. On a bounded range of scaled radii the true potential is at most
the reference quadratic plus $\epsilon k_p+O(1)$, with arbitrarily small fixed
$\epsilon>0$, by the frozen norm-density bound and static comparison.
Outside a sufficiently large bounded range, the zero-field projection
tail gives
\[
 \phi_p(y)-\phi_p(0)
 \le Ck_p\left(1+
       \left(\frac{\|y\|}{p^2M_p}\right)^{6/5}\right).
\]
This follows by integrating the tail
$\exp(-ak_pv^6)$ against the linear field tilt. Its growth is absorbed
by the reference quadratic. For the extension, use the support centered
at the point of radius $\ell_p(\alpha)$ in the direction of $y$.
At radial distance $vp^2M_p$ its excess above the reference is at most
\[
 \epsilon k_p+\epsilon k_pv+C(\eta+s_0)k_pv^2+O(1).
\]
Choose $\eta+s_0$ small compared with $\gamma^2$, and choose
$\epsilon$ after the fixed gap to the boundary of agreement. These errors
are strictly smaller than the Gaussian deviation cost outside agreement.

The normalizing integrals relative to the reference are bounded below
by $\exp(-\epsilon'k_p)$ with any needed fixed $\epsilon'>0$.
Indeed integrate over a sufficiently thin fixed relative shell around
$\ell_p(\alpha)$. Both potentials agree there, their values at the
central radius are given by static comparison, and their gradients are
bounded in units of $M_p$. Gaussian concentration gives a positive lower
bound for the reference mass of this shell for large $n$. This proves
the asserted transition tails; polynomial moments are absorbed by the
same Gaussian deviation estimate and the subquadratic bound above.
\end{proof}

The extension permits Brascamp--Lieb's covariance inequality
\cite{BrascampLieb1976} to be used in the Gaussian convolution. After
whitening, the transition has a fixed strict log-concavity margin depending
on $\gamma$, because
$\Delta\preceq(1-x^2)p^2\Id$ and $\eta\ll\gamma^2$. If $K$ denotes
the extended Hessian, this inequality bounds its covariance by the
expectation of
\[
 \Delta^{1/2}
 (\Id-\Delta^{1/2}K\Delta^{1/2})^{-1}\Delta^{1/2}.
\]
Mollification justifies this formula for the semiconcave extension.
Lemma~\ref{lem:variance-extension} then replaces the extended Hessian by
the true one inside agreement and replaces the extended transition by
the true transition, with exponentially small errors. The covariance
comparison includes the means: Cauchy--Schwarz and the polynomial moment
tail bound control the difference of first and second moments. In
\eqref{eq:source-18}, only $\Delta_c$ is inverted outside this expression,
so newly activated directions of arbitrarily small precision cause no loss.

\begin{lemma}[Removing the radial error]\label{lem:variance-radial}
On the coarse annulus, let $D=p^2K_p-\Id$, $e=h_p/\|h_p\|$,
$z=D_{\perp e}$, and define, on $e^\perp$,
\[
 E_p=(D_{\perp\perp}+zz^{\mathsf T}/|D_{ee}|)_+,
\]
extended by zero on $e$. Put $s_p^+=\|D_+\|$. Then
$D\preceq E_p$ and
\begin{equation}\label{eq:source-19}
 \|E_p\|\le Cs_p^+,
 \qquad
 \frac{\Tr E_p}{k_p}
 \le C\frac{\|D\|_{\mathrm F}}{\sqrt{k_p}}+\frac C{k_p}.
\end{equation}
\end{lemma}

\begin{proof}
The coarse radial deficit gives $D_{ee}\le-\varepsilon_1<0$.
For each transverse vector, maximize the quadratic form of $D$ over its
radial coordinate. The resulting form is
$D_{\perp\perp}+zz^{\mathsf T}/|D_{ee}|$, proving $D\preceq E_p$.
Since $K_p\succeq0$ and $p^2K_p\preceq(1+\varepsilon_0)\Id$, both
$D$ and the maximizing radial coordinate for a unit transverse vector
are bounded by constants. Comparing that vector with
$D\preceq s_p^+\Id$ proves the norm estimate. The trace of the rank-one
correction is bounded, while the trace of the positive part of
$D_{\perp\perp}$ is at most
$\sqrt{\dim\mathcal H_p}\,\|D\|_{\mathrm F}$. This proves
\eqref{eq:source-19}.
\end{proof}

Define $E_p=0$ outside the coarse annulus, and also whenever the local
covariance inequalities needed in Lemma~\ref{lem:variance-radial} fail.
This definition uses only the physical parent data $(W_p,h_p)$ and the
spectrum. In particular, it is defined and bounded before imposing
$\mathcal G_n$. On that event, the local failures never occur on the
annulus. By choosing $\varepsilon_0$ small enough, we have
$\|E_p\|\ll\gamma^2$.

Insert $K_p\preceq p^{-2}(\Id+E_p)$ in the covariance comparison.
The baseline $p^{-2}\Id$ gives $c^{-2}\Id$ in
\eqref{eq:source-18}. For the correction, expand the inverse around
$\Id-\Delta/p^2$. Its inverse has norm at most $x^{-2}$; because
$\|E_p\|\ll x^2$, the whole positive correction is at most twice the
first insertion in quadratic-form order. On the child subspace the two
outer inverse factors are scalars. Multiplication by $c^2$ yields
\begin{equation}\label{eq:source-20}
 s_c^+\le Cx^{-2}\mathsf T_{c,p,h_c}
                    [\|P_cE_pP_c\|]+C_\gamma e^{-a_1k_p}
\end{equation}
for every child field in its annulus. The first constant is absolute.

\begin{lemma}[Compression by the unused frame]\label{lem:variance-compression}
For $x=c/p\in[\gamma,2\gamma]$,
\begin{equation}\label{eq:source-21}
 \|P_cE_pP_c\|
 \le C\left(\frac{\Tr E_p}{k_p}+x^3\|E_p\|\right)
\end{equation}
outside a set of $\mathbf P$-probability at most
$C\exp(-a_\gamma k_p)$. This is an unconditional assertion with respect
to the unused spectral refinement; its exceptional set may afterwards
be intersected with $\mathcal G_n$.
\end{lemma}

\begin{proof}
Condition on $W_p$, the physical field $h_p$, and its spectral coordinates,
as in Lemma~\ref{lem:observations-angular}. The remaining frame in
$e^\perp\cap\mathcal H_p$ is Haar, and $E_p$ is fixed under this
conditioning. For a uniformly rotating unit vector $v$ in that space,
the Gaussian representation of the uniform sphere gives
\[
 \mathbb P\left(v^{\mathsf T}E_pv>
   \frac{C(\Tr E_p+t\|E_p\|)}{k_p}\right)
 \le C(e^{-t}+e^{-c_0k_p}).
\]
The same bound applies to a fixed vector of norm at most one. Choose a
constant-resolution net of the child unit sphere and project it onto
$e^\perp$. Its cardinality is at most
$\exp(C\dim\mathcal H_c)\le\exp(Cx^3k_p)$. With
$t=C_2x^3k_p$, a sufficiently large $C_2$ pays for the net; taking
$\gamma$ sufficiently small also pays for the second exponential term.
The net approximation for a positive quadratic form proves
\eqref{eq:source-21}.
\end{proof}

\begin{proof}[Proof of Proposition~\ref{prop:variance-curvature}]
Let $\widehat s_p$ be $s_p^+$ on the full coarse annulus and zero
outside it. We now restrict to $\mathcal G_n$. Average
\eqref{eq:source-20} over
$c/(2\gamma)\le p\le c/\gamma$ with normalized logarithmic measure,
and condition on $W$ and the observation path through $c$.
Combining \eqref{eq:source-19} and \eqref{eq:source-21}, the operator-norm
term has coefficient at most $Cx^{-2}x^3\le C\gamma$.
Choose $\gamma$ so that this coefficient, denoted by $\theta$, is less
than $1/2$. The trace term is controlled by Cauchy--Schwarz and the
integrated Frobenius estimate \eqref{eq:source-14}. We obtain
\begin{equation}\label{eq:source-22}
 \widehat s_c\le
 \theta\operatorname{avg}_{c/(2\gamma)\le p\le c/\gamma}
    \mathbf E[\widehat s_p\mid W,\mathcal F_c]
 +C_\gamma(c^{b}+k_c^{-b})+\operatorname{err}_c,
\end{equation}
where $b>0$, $\mathcal F_c$ is the observation filtration, and
$\operatorname{err}_c\ge0$ satisfies
\[
 \mathbf E[\mathbf1_{\mathcal G_n}\operatorname{err}_c]
 \le Ce^{-n^\xi}
\]
uniformly at deterministic scales. To verify the exceptional terms in
this assertion, first apply \eqref{eq:source-14} to the parent scale
integral before taking its conditional expectation. Its failures and
those of \eqref{eq:source-21} have exponential probability and at most
polynomial cost. Since $W$ is included in the conditioning, the indicator
of $\mathcal G_n$ commutes with that conditional expectation. No assertion
of Haar independence conditional on $\mathcal G_n$ is used. If the child
field is outside its annulus, the inequality holds by truncation.

Choose $a\le b$ sufficiently small that $\theta\gamma^{-a}<1$.
For a sufficiently large $C'$, the function
$b_0(c)=C'(c^a+k_c^{-a})$ dominates its parent average multiplied by
$\theta$, plus $C_\gamma(c^b+k_c^{-b})$, with a fixed margin.
It also dominates the coarse constant bound at the upper scales where
the parent interval no longer fits below $p_*$. Take positive parts
above $b_0$ in \eqref{eq:source-22}, use conditional Jensen, and average
over $\mathbf P$ with $\mathbf1_{\mathcal G_n}$. At each iteration the
scale increases by at least $(2\gamma)^{-1}>1$, so there are only
$O(\log n)$ generations; their errors sum to at most
$Ce^{-n^\xi}\sum_{j\ge0}\theta^j$. Thus
\[
 \mathbf E[\mathbf1_{\mathcal G_n}
       (\widehat s_c-b_0(c))_+]\le Ce^{-n^\xi}.
\]
Enlarging $C'$ and applying Markov's inequality proves the asserted
probability bound, since $b_0(c)$ is bounded below by an inverse
polynomial on the scale range. Lemma~\ref{lem:observations-annulus}
removes the truncation. All bounds were uniform at deterministic scales,
so Fubini's theorem gives their integrated version.
\end{proof}

\subsection{The variance bound at a fixed upper scale}

We next obtain the starting variance estimate from the fixed-scale
continuation in \cite[Section~19.4]{CM}. We record precisely why its
entropy estimate applies to the perturbations needed here.

\begin{lemma}[The fixed-scale starting bound]\label{lem:variance-start}
On disorder events of probability tending to one, for a sufficiently
small fixed $p_{\mathrm{top}}>0$ and every real function $f$ on the cube,
\[
 V_{p_{\mathrm{top}}}(f)
 \le C\cE(f)+Ce^{-a n}\|f\|_\infty^2.
\]
The event and constants are independent of $f$.
\end{lemma}

\begin{proof}
The continuation in \cite[Section~19.4]{CM} increases the observation
precision from $B_{p_{\mathrm{top}}}$ through finitely many cap intervals,
a bridge to a classical field, and a bounded-duration continuation to
a diagonal-interaction endpoint. For an initial change of spin law from
$\mu$ to $\nu$, let $\nu_s,\mu_s$ be the two posteriors, let
$l_s=D(\nu_s\Vert\mu_s)$, and put
$H(s)=E_{Q_\nu}l_s$. On this continuation, the estimate used is
\begin{equation}\label{eq:variance-entropy-ode}
 \dot H(s)\ge-CH(s)-e^{-a n}.
\end{equation}
This is \cite[Equation~(19.17)]{CM}. Its proof uses upper bounds on
observation likelihoods and posterior entropy tails, rather than a
support restriction on the posterior. We explain its application when
the initial density lies between $1/2$ and $2$. Every observation likelihood
then has the same bounds, every posterior density ratio lies between
$1/4$ and $4$, and $l_s\le\log4$.

At the classical part, the ordinary posterior covariance is bounded on
a fixed normalized field neighborhood outside an event of probability
$e^{-an}$; farther along the continuation it is bounded for every field.
The linear entropy inequality in that neighborhood gives
$\|m_\nu-m_\mu\|^2\le Cl_s$ whenever $l_s$ is bounded: a trial field
shift of size proportional to $\sqrt{l_s}$ fits the neighborhood.
The bounded density ratios pay for the ordinary exceptional paths.

To propagate this mean estimate across a preceding cap or bridge
interval, use the approximate gradient $g_s$ furnished by the extended
potential there. The transport identity
\cite[Lemma~18.3]{CM} holds under either initial law. If $U$ is its
linear transport from the current time $s$ to the end of the interval,
subtraction gives, for a vector $v$ in the current rate subspace,
\begin{align*}
 \langle v,m_{\nu,s}-m_{\mu,s}\rangle
 ={}&E_{Q_{\nu,h}^{\mathrm{cont}}}
       \langle Uv,m_{\nu,\mathrm{end}}-m_{\mu,\mathrm{end}}\rangle\\
 &+(E_{Q_{\nu,h}^{\mathrm{cont}}}
       -E_{Q_{\mu,h}^{\mathrm{cont}}})
       \langle Uv,m_{\mu,\mathrm{end}}-g_{\mathrm{end}}\rangle.
\end{align*}
Here the two continuation laws live on the same path space, and the same
path function $U$ is used in both expectations. Its operator norm is
bounded by a fixed constant on these fixed intervals. At agreement,
the terminal approximation error in the tested subspace has bounded
norm. Truncate it there. Data processing gives
$D(Q_{\nu,h}^{\mathrm{cont}}\Vert Q_{\mu,h}^{\mathrm{cont}})\le l_s$,
so Pinsker's inequality bounds the second term by
$C\sqrt{l_s}\|v\|$.

The complements have exponentially small averaged squared cost. Indeed,
for any nonnegative terminal cost $W_1$, the bounded full-path and
prefix likelihoods give both
\[
 E_{Q_\nu}E_{Q_{\nu,h}^{\mathrm{cont}}}W_1
 \le C E_{Q_\mu}W_1,
 \qquad
 E_{Q_\nu}E_{Q_{\mu,h}^{\mathrm{cont}}}W_1
 \le C E_{Q_\mu}W_1.
\]
Apply these inequalities to the polynomially weighted ordinary
exceptions in the terminal approximation estimates of
\cite[Section~19.4]{CM}. Conditional Cauchy--Schwarz controls the first
transport term by the terminal squared mean difference. Taking the
supremum over unit $v$ costs no dimension factor because all the bounds
are norm bounds uniform in $v$. Since terminal expected posterior entropy
is at most $H(s)$, backward induction through the finite list yields
\[
 E_{Q_\nu}\|P_{\mathrm{rate}}
                   (m_{\nu,s}-m_{\mu,s})\|^2
 \le CH(s)+Ce^{-a n}.
\]
The observation entropy identity
\cite[Lemma~19.2]{CM}, with bounded precision rates on these intervals,
now proves \eqref{eq:variance-entropy-ode}. This also proves that its
constants apply uniformly to the bounded-density perturbations below.

Normalize $\|f\|_\infty\le1$, set $\eta=e^{-n^{1/3}}$, and take
\[
 \frac{d\nu}{d\mu}=\frac{1+\eta f}{1+\eta\mu f}.
\]
Every path likelihood and every posterior density ratio is bounded by
fixed constants. Taylor's formula for relative entropy gives, uniformly
at every posterior,
\[
 c\eta^2\Var_{\mu_s}(f)
 \le D(\nu_s\Vert\mu_s)
 \le C\eta^2\Var_{\mu_s}(f).
\]
Thus $H(s)$ is comparable to $\eta^2V_s(f)$. Integrating
\eqref{eq:variance-entropy-ode} over a bounded duration and dividing by
$\eta^2$ gives
$V_{p_{\mathrm{top}}}(f)\le CV_{\mathrm{end}}(f)+Ce^{-a'n}$.
At the endpoint the posterior is a product measure, so variance
tensorization and conditional-variance contraction give
\[
 V_{\mathrm{end}}(f)
 \le E\sum_i\Var(f\mid X_{-i},\text{endpoint observations})
 \le\sum_i\mu[\Var(f\mid X_{-i})]=\cE(f).
\]
Rescaling $f$ proves the lemma.
\end{proof}

\subsection{Integration of the posterior variance identity}

The passage from a posterior covariance bound to a variance bound uses
the localization mechanism of
\cite[Lemma~9, arXiv version~2]{EldanKoehlerZeitouni2022}.
For the deterministic precision path considered here, it takes the following
form.

\begin{proof}[Proof of Proposition~\ref{prop:variance-main}]
For fixed $W$, the posterior mean of $f$ is a martingale in the
observation filtration. Its diffusion coefficient, with parameter
$t=p^2$, is $P_p\Cov(X,f\mid h_p)$. Consequently
\[
 -\frac{d}{d(p^2)}V_p(f)
 =Q_\mu\|P_p\Cov(X,f\mid h_p)\|^2
 \le Q_\mu\bigl[\|K_p\|\Var(f\mid h_p)\bigr].
\]
Integrate the exceptional probabilities in
Proposition~\ref{prop:variance-curvature} in logarithmic scale, then apply
Markov's inequality and Lemma~\ref{lem:edge-biased-law}. On disorder
events of probability arbitrarily close to one, intersected with
$\mathcal G_n$, their total ordinary path mass is at most
$Ce^{-n^\xi}$. The same statement holds for the annulus exceptions
used with the coarse bound between $p_*$ and $p_{\mathrm{top}}$.
These events are chosen before, and independently of, $f$.

On the exceptional paths the elementary estimate
$\|P_p\Cov(X,f\mid h_p)\|^2\le Cn\|f\|_\infty^2$ suffices.
On the remaining paths, put $\epsilon_p=C(p^a+k_p^{-a})$ below $p_*$
and use the fixed coarse tolerance above it. Backward integration from
$p_{\mathrm{top}}$ gives an amplification bounded by
\[
 \exp\left(\int_{p^2}^{p_{\mathrm{top}}^2}
                  \frac{1+\epsilon_{\sqrt t}}{t}\,dt\right)
 \le Cp^{-2},
\]
because
\[
 \int_r^{p_*}(p^a+k_p^{-a})\,d\log p\le C,
\]
and the remaining scales form a fixed interval bounded away from zero.
The amplification and the elementary exceptional-path costs are
polynomial in $n$; they are absorbed by decreasing $\xi$.
Lemma~\ref{lem:variance-start} therefore gives
\eqref{eq:source-16}. The same integrated exceptional event works for
every lower endpoint $p$, so the assertion is simultaneous in scale as
well as in the test.
\end{proof}
